% Written from the committed results in data/llm_defender/claude-sonnet-5 (numbers are generated macros).
\paragraph{Results} Table~\ref{tab:llm}. The run used \ClEpisodes{} episodes: a spending
cap of about eight US dollars stopped it before the last two repeats of two uncertifiable
regimes (total cost \$\ClCost{}). Against the typical adversary Claude certified
\ClTypCert{} episodes, with \ClTypMean{} mitigations on average against \ClTypGreedyMean{}
for the greedy defender, \ClTypCoverageMean{} for the coverage heuristic and
\ClTypOptMean{} for the exact optimum (\ClTypRatio{} times the optimum per episode). Its
stated reasons show a coverage-first strategy: it adds the mitigations covering the most
techniques, and its portfolios resemble the coverage heuristic's. Against the adaptive
adversary the picture reverses. In the one regime that admits a certificate it certified
\ClAdaptCert{} episodes with \ClAdaptRange{} mitigations, where the greedy and coverage
heuristics need \ClAdaptGreedy{} and the exact optimum is \ClAdaptOpt. In the regimes no
portfolio can certify it never claimed success (it cannot: the score is a certificate),
and it chose to stop in \ClUncertStop{} episodes, but only after \ClUncertSteps{} turns on
average and with \ClUncertDeployed{} mitigations deployed, so it recognised impossibility
late. It made \ClInvalidCount{} invalid action in \ClSteps{} turns. With three repeats per
regime these are descriptive results about one model, not estimates for language-model
defenders in general.
